%%%%% slides template
% Giovanni Ramirez (ramirez@ecfm.usac.edu.gt)
% 20200828
% Require: images/ecfmByN.png images/usacByN.jpg
%
%
% Use this option to print notes only
% \documentclass[xcolor=dvipsnames,handout,notes=only]{beamer}%
%
% Use this option to print notes and slides
% \documentclass[xcolor=dvipsnames,handout,notes=show]{beamer}%
%
% Use this option to get a printer-friendly version
% \documentclass[xcolor=dvipsnames,handout]{beamer}%
%
% Unse this option to get slides only
\documentclass[xcolor=dvipsnames,presentation]{beamer}%
%
%%%%%


%%%%% theme and coloring
% a list of themes can be found here https://hartwork.org/beamer-theme-matrix/
% I like a simple Malmoe with some modifications and color-schemes like
% dolphin for a white background, beetle for a gray background or dove for
% plain slides
%
% theme definition
\usetheme{Malmoe}
%
% definition of the presentation mode: full slides
\mode<presentation> {
% coloring for slides other options: beetle, dove
  \usecolortheme{dolphin}
% to include a slide with the table of contents
 \AtBeginSection[] {
  \begin{frame}
    \tableofcontents[currentsection]
  \end{frame} }
}
%
% definition of the handout mode: printer-friendly
\mode<handout>{
% coloring: dove is white
  \usecolortheme{dove}
% this is to print 2 slides in 1 page, other values are allowed
  \usepackage{pgfpages}
  \renewcommand\textbullet{\ensuremath{\bullet}}
  \pgfpagesuselayout{2 on 1}[letterpaper,border shrink=5mm]
% set the fontsize of the notes
  \setbeamerfont{note page}{size=\scriptsize}
% disable the slide with the table of contents
  \AtBeginSection[]{}
}
%%%%%

%%%%% personalisation
% clear all headers
\setbeamertemplate{headline}{}
% disable navigation buttons
\setbeamertemplate{navigation symbols}{}
% configure the blocks used for equations, theorems, etc
\setbeamertemplate{blocks}[rounded][shadow=true]
% clear the footline
\setbeamertemplate{footline}{}
% set the footline
\setbeamertemplate{footline}{
  \hbox{%
% left block with the shortname, it is set in the \author command
    \begin{beamercolorbox}[wd=.40\paperwidth,ht=2.25ex,dp=1ex,center]%
      {author in head/foot}%
      \usebeamerfont{author in head/foot}\insertshortauthor
    \end{beamercolorbox}%
% center block with the short title, it is set in the \title command
    \begin{beamercolorbox}[wd=.50\paperwidth,ht=2.25ex,dp=1ex,center]%
      {title in head/foot}%
      \usebeamerfont{title in head/foot}\insertshorttitle
    \end{beamercolorbox}%
% right block with the number of slides
    \begin{beamercolorbox}[wd=.10\paperwidth,ht=2.25ex,dp=1ex,right]%
      {date in head/foot}%
      \usebeamerfont{date in head/foot}
      \insertframenumber{} / \inserttotalframenumber\hspace*{1em}
    \end{beamercolorbox}}%
  \vskip0pt%
}
%%%%%

%%%%% packages
%
% language, decimal point mark
\usepackage[spanish,es-nodecimaldot]{babel}%
%
% to use only T1 scalable fonts
\usepackage[T1]{fontenc}%
\let\Tiny=\tiny% this is a proper configuration of the tiny font size
%
% to use UTF coding, other options latin1
\usepackage[utf8]{inputenc}%
%
\usepackage{array}
% special math symbols
\usepackage{latexsym,amsfonts,amsmath}%
%
% to include graphics
\usepackage{graphicx}%
% path to images
\graphicspath{ {../images/} }
%
% to use tiks: nice to draw diagrams
\usepackage{tikz}%
\usetikzlibrary{arrows,shapes,positioning}
%
% to use SI units and physics notation
\usepackage{physics}%
\usepackage{multirow}
\usepackage{xcolor}
\usepackage{subcaption}
%
% to use hypenation
\usepackage{ragged2e}%
\let\raggedright=\RaggedRight%
%
\usepackage{dsfont}

\newcommand{\subind}[2]{{}_{#1}^{#2}}
%%%%% Talk data
%
%title
\title[Canales cuánticos PCE]% this is the short title
{Mapeos proyectivos en sistemas de varios qubits}% this is the long title
\subtitle{Operaciones PCE \newline\vspace{1.4cm}}
%
% author
\author[José Alfredo de León]% this is the short name
{José Alfredo de León}% this is the long name
%
% affiliation
\institute{Asesorado por:\\Dr. Carlos Francisco Pineda Zorrilla (IFUNAM)
\newline y M.Sc. Juan Diego Chang (ECFM-USAC)}%
%
% date and information about the conference
\date{XX octubre 2021}
%%%%%


%%%%%
\makeatletter
\newenvironment{myitemize}{%
   \setlength{\topsep}{0pt}
   \setlength{\partopsep}{0pt}
   \renewcommand*{\@listi}{\leftmargin\leftmargini \parsep\z@ \topsep\z@ \itemsep\z@}
   \let\@listI\@listi
   \itemize
}{\enditemize}
\makeatother 


\begin{document}
% % % % % % % % % % % % % % % % % % % % % % % % % % % % % % % % % % % % % %
% Title and Outline
% % % % % % % % % % % % % % % % % % % % % % % % % % % % % % % % % % % % % %
\begin{frame}[plain]
  \titlepage 
  \begin{tikzpicture}[x=1mm,y=1mm,overlay,remember picture]
    \pgftransformshift{\pgfpointanchor{current page}{center}}
    \node[inner sep=0pt] (usac) at (-40,-14.5) %
    {\includegraphics[height=14mm]{logos/ecfmByN}};
    \node[inner sep=0pt] (usac) at (39,-14.5) %
    {\includegraphics[height=14mm]{logos/ifunam}};
    \node[inner sep=0pt] (usac) at (0,11) %
    {\includegraphics[height=14mm]{portada}};
  \end{tikzpicture}
\end{frame}

\begin{frame}{Agradecimientos}
	\begin{itemize}
		\item A mis padres.
		\item A mis asesores: Carlos y Juan Diego. 
		\item A Alejandro Fonseca y David Dávalos.
		\item A mis prefesores de la ECFM.
		\item A Cindy.
		\item A mis mejores amigos Benja, Gómez, Papaya y José Guillermo.
		\item A mis amigos de la U.
	\end{itemize}
\end{frame}

\begin{frame}{Esquema de la presentación}
  \tableofcontents[hidesubsections] 
\end{frame}


% % % % % % % % % % % % % % % % % % % % % % % % % % % % % % % % % % % % % %
% Introduction and Motivation
% % % % % % % % % % % % % % % % % % % % % % % % % % % % % % % % % % % % % %
\section{Introducción}
\label{sec:Intro}
\begin{frame}{Motivación}
	\begin{itemize}
	\item En una descripción más completa de un sistema cuántico se 
	debe de considerar con su entorno
	\item La decoherencia es un problema interesante
	\item Operaciones PCE 
	\end{itemize}
%	\only<1>{
%	Sistema ideal:
%	\vspace{-.5cm}
%	\begin{columns}
%	\begin{column}{.5\textwidth}
%	\begin{center}
%		\includegraphics[width=6cm]{sistema_cerrado.pdf}
%	\end{center}
%	\end{column}
%	\hfill
%	\begin{column}{.45\textwidth}
%	\begin{center}
%		\includegraphics[width=5cm]{dynamics_closed}
%	\end{center}
%	\end{column}
%	\end{columns}
%	}
%\only<2>{
%	Sistemas reales: abiertos
%	\vspace{-.1cm}
%	\begin{center}
%		\includegraphics[width=5cm]{sistema_abierto01.pdf}
%		\hfill
%		\includegraphics[width=5cm]{sistema_abierto02.pdf}
%	\end{center}
%	\begin{center}
%		\hspace{.5cm}
%		\includegraphics[width=5cm]{dynamics_open}
%	\end{center}
%  }

\end{frame}


\section{Fundamentos teóricos}
\begin{frame}{Matriz de densidad}
%{¿cómo representar a los estados de un sistema abierto?}
	\vfill
%   \only<1-2>{
   \begin{itemize}
   	\item  La matriz de densidad $\rho$ es la herramienta apropiada para 
   	describir al estado un sistema abierto. \vfill
   	\item La matriz $\rho$ es una matriz de densidad si y sólo si
	\begin{enumerate}
		\item $\Tr (\rho)=1$,
%		\item $\rho^{\dagger}=\rho$,
		\item $\matrixel{\psi}{\rho}{\psi}\geq 0$.
	\end{enumerate}	\vfill
%	}
%		\onslide<2>{
	\item 
	Por ejemplo, si un sistema se encuentra en el estado $\ket{\psi}$,
  \begin{align*}
  \rho=\dyad{\psi}{\psi}.
  \end{align*}\vfill
%	}
   \end{itemize}	
\end{frame}

\begin{frame}{Qubits}
	\begin{columns}	\hspace{.4cm}
	\begin{column}{0.6\textwidth}
	\begin{itemize}
		\item Un qubit es un sistema cuántico de dos niveles,
		\begin{align*}
			\rho &= \frac{\mathds{1}+r_1\sigma_x+r_2\sigma_y+r_3\sigma_z}{2}.
		\end{align*}	
		\item $n$ qubits:
		\begin{align*}
		\rho = \frac{1}{2^n}\sum _{j_1,\ldots,j_n=0}^3 r_{j_1,\ldots,j_n}
		\sigma_{j_1}\otimes \ldots
		\otimes\sigma_{j_n},%\qquad r_{0,\ldots,0}=1,
		\end{align*}
		$r_{j_1,\ldots,j_n}$ componentes de Pauli.
%		\item $\qty(r_1,r_2,r_3)$ especifican las coordenadas cartesianas de un 
% 		punto en la esfera de Bloch, $\alert{\rho = \qty(1,r_1,r_2,r_3)}$
	\end{itemize} \vfill
	\end{column}%\hspace{-1cm}
	\begin{column}{0.4\textwidth}  
  		\begin{center}
    		\includegraphics[width=.85\textwidth]{bloch_sphere}      
     \end{center}
     \vspace{3cm}
		\end{column}
	\end{columns}
%	\begin{itemize}
%		\item Para $n$ qubits
%		\begin{align*}
%		\rho = \frac{1}{2^n}\sum _{j_1,\ldots,j_n=0}^3 r_{j_1,\ldots,j_n}
%		\sigma_{j_1}\otimes \ldots
%		\otimes\sigma_{j_n},\qquad r_{0,\ldots,0}=1,
%		\end{align*}
%		$r_{j_1,\ldots,j_n}$ componentes de Pauli.
%	\end{itemize}
\end{frame}

\begin{frame}{Canales cuánticos}%
  %{¿cómo describir la dinámica de los sistemas abiertos?}
  \vfill
  \begin{itemize}
  	\item La teoría de los canales cuánticos es un formalismo para describir 
  la evolución de los sistemas abiertos,
  \begin{align*}
	\mathcal{E}(\rho)=\rho'.
	\end{align*} \vfill
  \only<1>{\item Una operación lineal $\mathcal{E}$ es un canal cuántico 
  si y sólo si
	\begin{enumerate}
	\item Preserva las características de la matriz de densidad
	\item Es una operación completamente positiva,
	$$
	\qty(\mathcal{E}\otimes \mathds{1})\qty[\dyad{\text{Bell}}{\text{Bell}}]\geq 0.
	$$
	\end{enumerate} \vfill
	} 
  \end{itemize}

\end{frame}	


\begin{frame}{Canales cuánticos de 1 qubit}
	Rotaciones de la esfera de Bloch
	\begin{center}
	\begin{tabular}{m{2.5cm} m{1.8cm} m{2.5cm}}
		\includegraphics[width=2.5cm]{bloch_sphere2}
		& \hspace{\fill} \Large$\overset{\normalsize{R_x(\pi/4)}}{\longmapsto}$ \hspace{\fill}
		& \includegraphics[width=2.5cm]{bloch_sphere_Rx_pi_medios2}
	\end{tabular}
	\end{center}
	
	Canal \textit{bit-flip}	 $\mathcal{E}$
	\begin{center}
	\begin{tabular}{m{2.5cm} m{1.8cm} m{2.5cm}}
		\includegraphics[width=2.5cm]{unit_sph}
		& \hspace{\fill} \Large{$\overset{\mathcal{E}}{\longmapsto}$} \hspace{\fill}
		& \includegraphics[width=2.5cm]{bit_flip_p0_3}
	\end{tabular}
	\end{center}
\end{frame}


%\begin{frame}{Completa positividad}{¿Para qué o qué?}
%	Supongamos una operación $\mathcal{E}$ que actúa sobre 
%	el sistema A.
%	
%	\only<1->{
%	\begin{center}
%	\includegraphics[height=6cm]{CP2}
%	\end{center}	
%	}
%	
%	No es suficiente que 	$\mathcal{E}(\rho_A)\geq0$, \alert{también 
%	$\qty(\mathcal{E}\otimes \mathds{1})\qty[\dyad{\text{Bell}}{\text{Bell}}]$.}
%\end{frame}
%
%\begin{frame}{Canal cuántico}
%	\begin{center}
%		\includegraphics[width=7cm]{QuantumChannel}
%	\end{center}
%	\vspace{-.2cm}
%	\only<1>{
%	Una operación $\mathcal{E}$ es un canal cuántico si y sólo si 
%	\begin{align*}
%	 \qty(\mathcal{E}\otimes \mathds{1})\qty[\dyad{\text{Bell}}{\text{Bell}}]
%	\hspace{.5cm}	 
%	 \alert{\text{(}\propto  \text{matriz de Choi de }\mathcal{E})}
%	\end{align*}
%	es (1) de traza unitaria, (2) Hermítica y (3) positiva.	
%	}
%\end{frame}
%
%
%\begin{frame}{Resumen}{Sistemas cuánticos abiertos}
%	\begin{itemize}
%		\item Estados cuánticos: matriz de densidad $\rho$
%		\vspace{1cm}
%		\item Evolución: canales cuánticos (operaciones CPTP)
%	\end{itemize}
%	\vfill
%\end{frame}

% % % % % % % % % % % % % % % % % % % % % % % % % % % % % % % % % % % % % %
% Qubits
% % % % % % % % % % % % % % % % % % % % % % % % % % % % % % % % % % % % % %

\section{Operaciones PCE}

%\begin{frame}{Canal \textit{bit-flip}}
%	El \textit{bit-flip} es otro canal cuántico de 1 qubit que actúa sobre 
%	la esfera de Bloch como
%	\begin{center}
%	\begin{tabular}{m{2.7cm} m{1.8cm} m{3cm}}
%		\includegraphics[width=3cm]{unit_sph}
%		& \hspace{\fill} \LARGE{$\overset{\mathcal{E}}{\longmapsto}$} \hspace{\fill}
%		& \includegraphics[width=3cm]{bit_flip_p0_3}
%	\end{tabular}
%	\end{center}
%%	
%%	Transforma a las componentes $r_i$ de la matriz de densidad como
%	\begin{align*}
%	\mathcal{E}:\qty(1,r_1,r_2,r_3) \longmapsto \qty(1,r_1,\qty(1-2p)r_2,\qty(1-2p)r_3)
%	\end{align*}
%\end{frame}

%\begin{frame}{Matriz de densidad de $n$ qubits}
%	En la base de productos tensoriales de las matrices de Pauli, 
%	la matriz de densidad de $n$ qubits se escribe
%	\begin{align*}
%		\rho = \frac{1}{2^n}\sum _{j_1,\ldots,j_n=0}^3 r_{j_1,\ldots,j_n}
%		\sigma_{j_1}\otimes \ldots
%		\otimes\sigma_{j_n},\qquad r_{0,\ldots,0}=1,
%	\end{align*}
%	\alert{$r_{j_1,\ldots,j_n}$: ``componentes de Pauli''.}
%\end{frame}

% % % % % % % % % % % % % % % % % % % % % % % % % % % % % % %
% Operaciones PCE
% % % % % % % % % % % % % % % % % % % % % % % % % % % % % % % 
%\section{Operaciones PCE}
%\label{sec:Theory}

%\begin{frame}{Motivación}{Operaciones PCE}
%\only<1>{
%	Un caso particular del canal \textit{bit-flip} es cuando
%	\begin{center}
%	\begin{tabular}{m{2.5cm} m{1.5cm} m{2.5cm}}
%		\includegraphics[width=3cm]{unit_sph}
%		& \hfill \LARGE{$\longmapsto$}
%		& \includegraphics[width=3cm]{bit_flip_p0_5}
%	\end{tabular}
%	\end{center}
%	
%	Las componentes de Pauli $r_i$ se transforman como
%	\begin{align*}
%	\qty(1,r_1,r_2,r_3)\longmapsto \qty(1,r_1,0,0).
%	\end{align*}
%	
%	\alert{¿Son canales cuánticos todas las operaciones que borran
%	cualesquiera de las componentes de Pauli de 1 qubit?}
%}
%
%\only<2>{
%	\alert{No.}
%	Las operaciones que borran dos componentes de Pauli no son canales 
%	cuánticos.
%	\begin{center}
%	\begin{tabular}{m{2.5cm} m{1.3cm} m{2.5cm}}
%		\includegraphics[width=3cm]{unit_sph}
%		& \hfill \LARGE{$\longmapsto$}
%		& \includegraphics[width=3.5cm]{unit_disk_xy}
%	\end{tabular}
%	\end{center}
%	\vspace{-1cm}
%	\begin{align*}
%	\qty(1,r_1,r_2,r_3)\longmapsto\qty(1,r_1,r_2,0).
%	\end{align*}
%	
%%	\alert{$\qty(\Lambda\otimes \mathds{1})\qty[\dyad{\text{Bell}}{\text{Bell}}]\ngeq0$, 
%%	por consiguiente	$\Lambda$ no es completamente positiva.}
%	\alert{No son completamente positivas.}
%		
%	\vspace{.5cm}
%}´
%\end{frame}


\begin{frame}{Operaciones PCE}
	\begin{itemize}
		\item Una operación PCE (\textit{Pauli component erasing}) 
		es una operación lineal que transforma a las 	componentes 
		de Pauli de una matriz de densidad $\rho$ de $n$ qubits como
		\begin{align*}
		r_{j_1,\ldots,j_n}\longmapsto \tau_{j_1,\ldots,j_n}r_{j_1,\ldots,j_n},
		\qquad \tau_{j_1,\ldots,j_n}=0,1.
		\end{align*} 
	
		\item Figuras PCE:
		\begin{figure} % {{{
		\begin{tikzpicture}[x=0.5cm, y=0.5cm] % {{{
		\pgfmathsetmacro{\unitstep}{5.8}
		\begin{scope}[shift={(0*\unitstep,0)}] % Coordenadas {{{
		\node at (-0.6,0.5) {(a)} ;
		    \foreach \x in {0,1,2,3} {
		      \foreach \y in {0,1,2,3} {
		        \begin{scope}[shift={(\x,-\y)}] 
		          \draw (0,0) rectangle (1,1); 
		          \node at (0.5,0.5) {\scalebox{.6}{$\tau_{\y,\x}$}};
		         \end{scope}
		%         \node at (0,-\y) (input\y) {$i_\y$};
		%         \node[block] at (2,-\y) (block\y) {$f_\y$};
		%         \draw[->] (input\y) -- (block\y);
		%         \draw[->] (block\y.east) -- +(0.5,0);
		    }
		    }
		\end{scope} % }}}
		\begin{scope}[shift={(1*\unitstep,0)}] % Good channel {{{
		\node at (-0.6,0.5) {(b)} ;
		    \foreach \x in {0,1,2,3} {
		      \foreach \y in {0,1,2,3} {
		        \begin{scope}[shift={(\x,-\y)}] 
		          \draw (0,0) rectangle (1,1); 
		%           \node at (0.5,0.5) {$\tau_{\y,\x}$};
		         \end{scope}
		%         \node at (0,-\y) (input\y) {$i_\y$};
		%         \node[block] at (2,-\y) (block\y) {$f_\y$};
		%         \draw[->] (input\y) -- (block\y);
		%         \draw[->] (block\y.east) -- +(0.5,0);
		    }
		    }
		\begin{scope}[shift={(0,-3)}] \fill[black] (0,0) rectangle (1,1); \end{scope}
		\begin{scope}[shift={(3,-3)}] \fill[black] (0,0) rectangle (1,1); \end{scope}
		\begin{scope}[shift={(0,0)}] \fill[black] (0,0) rectangle (1,1); \end{scope}
		\begin{scope}[shift={(3,0)}] \fill[black] (0,0) rectangle (1,1); \end{scope}
		\end{scope} % }}}
		\begin{scope}[shift={(2*\unitstep,0)}] % Good channel {{{
		\node at (-0.6,0.5) {(c)} ;
		    \foreach \x in {0,1,2,3} {
		      \foreach \y in {0,1,2,3} {
		        \begin{scope}[shift={(\x,-\y)}] 
		          \draw (0,0) rectangle (1,1); 
		%           \node at (0.5,0.5) {$\tau_{\y,\x}$};
		         \end{scope}
		%         \node at (0,-\y) (input\y) {$i_\y$};
		%         \node[block] at (2,-\y) (block\y) {$f_\y$};
		%         \draw[->] (input\y) -- (block\y);
		%         \draw[->] (block\y.east) -- +(0.5,0);
		    }
		    }
		\begin{scope}[shift={(0,0)}] \fill[black] (0,0) rectangle (1,1); \end{scope}
		\begin{scope}[shift={(0,-1)}] \fill[black] (0,0) rectangle (1,1); \end{scope}
		\begin{scope}[shift={(2,0)}] \fill[black] (0,0) rectangle (1,1); \end{scope}
		\begin{scope}[shift={(2,-2)}] \fill[black] (0,0) rectangle (1,1); \end{scope}
		\begin{scope}[shift={(2,-3)}] \fill[black] (0,0) rectangle (1,1); \end{scope}
		\end{scope} % }}}
		\end{tikzpicture} % }}}
		\end{figure} % }}}

		\item El problema de las operaciones PCE...
	\end{itemize}
\end{frame}


\begin{frame}[t]{Método numérico}
%	\begin{itemize}
%		\item Contar que se desarrolló un paquete en Mathematica para implementar 
%		el método numérico y un .nb para probar y reproducir nuestros resultados.
%		Esto sería sin entrar en detalles, sólo tipo "se hizo este paquete con un montón 
%		de rutinas y se hizo un notebook para que cualquiera pueda reproducir 
%		nuestros resultados y usar nuestras rutinas".
%		Esta diapositiva es más con el objetivo de mostrarle a los evaluadores 
%		el trabajo que hice.
%	\end{itemize}
	\only<1>{
	\begin{columns}
		\begin{column}{0.5\textwidth}
			\begin{figure}
		\centering
		\includegraphics[width=1\textwidth]{numerico_01}
	\end{figure}		
		\end{column}
		\begin{column}{0.5\textwidth}
			\begin{figure}
		\centering
		\includegraphics[width=1\textwidth]{numerico_02}
	\end{figure}		
		\end{column}
	\end{columns}		
	}
	\only<2>{
	https://github.com/deleonja/projective\_maps
	\begin{figure}
		\centering
		\includegraphics[width=.9\textwidth]{numerico_03}
	\end{figure}
	}
	
\end{frame}



\begin{frame}[t]{Resultados}
	Regla $2^k$: los canales cuánticos PCE preservan una cantidad 
		de componentes de Pauli que son potencias de dos.
	\begin{columns}[t]
		\begin{column}{0.3\textwidth}
			\begin{itemize}
				\item 1 qubit:
			\end{itemize}
			\only<1>{
			\begin{figure}
				\includegraphics[height=1.4cm]{1qubit}
			\end{figure}
			}
			\only<2>{
				\begin{figure}
					\centering
					\includegraphics[height=1.15cm]{id_1qubit}
				\end{figure}\vspace*{-.5cm}\hspace*{.75cm}\rule{1.8cm}{1.5pt}\vspace*{-.2cm}
				\begin{figure}
					\centering
					\includegraphics[height=1.4cm]{other_1qubit}
				\end{figure}\vspace*{-.65cm}\hspace*{.75cm}\rule{1.8cm}{1.5pt}\vspace*{-.2cm}
				\begin{figure}
					\centering
					\includegraphics[height=1.15cm]{dep_1qubit}
				\end{figure}
			}				
		\end{column} 
		\begin{column}{0.7\textwidth}		
		\begin{itemize}
		\item 2 qubits:
		\end{itemize}
		\only<1>{
		\begin{figure}
			\centering
			\includegraphics[width=.8\textwidth]{muestra}
		\end{figure}
		}
%		\vspace{-1.1cm}
%		\begin{figure}
%		\centering
%			\begin{minipage}{0.5\textwidth}
%				\centering
%				\hspace{.5cm}
%				\includegraphics[width=.6cm]{C1}
%		%		\subcaption{C$\subind{1}{}$}
%			\end{minipage}\\ \noindent\rule{1.05\textwidth}{0.3pt} \\ \vspace{.1cm} 
%			\begin{minipage}{0.5\textwidth}
%				\includegraphics[width=3.8cm]{C2_1}
%				\vspace{.6cm}
%		%		\subcaption{C$\subind{2}{1}$}
%			\end{minipage} 
%			\begin{minipage}{0.4\textwidth}
%				\includegraphics[width=3.8cm]{C2_2}
%		%		\subcaption{C$\subind{2}{2}$}
%			\end{minipage} \\ \noindent\rule{1.05\textwidth}{0.3pt} \\ \vspace{.1cm} 
		\only<2>{
		\begin{figure}
			\centering
			\includegraphics[height=.81cm]{C4_1}
		\end{figure}\vspace*{-.75cm}\hspace*{1cm}\rule{5.5cm}{1.5pt}\vspace*{-.2cm}
		\begin{figure}
			\centering
			\includegraphics[height=1.4cm]{C4_2}
		\end{figure}\vspace*{-.65cm}\hspace*{1cm}\rule{5.5cm}{1.5pt}\vspace*{-.2cm}
		\begin{figure}
			\centering
			\includegraphics[height=2.1cm]{C4_3}
		\end{figure}\vspace*{-.65cm}\hspace*{1cm}\rule{5.5cm}{1.5pt}\vspace*{-.2cm}
		\begin{figure}
			\centering
			\includegraphics[height=.81cm]{C4_4}
		\end{figure}
%			\begin{minipage}{0.5\textwidth}
%				\centering
%				\includegraphics[height=1cm]{C4_1}
%				\vspace{.6cm}
%		%		\subcaption{C$\subind{4}{1}$}
%			\end{minipage} 
%			\begin{minipage}[3.1cm]{0.4\textwidth}
%				\includegraphics[width=.45\textwidth]{C4_2}
%		%		\subcaption{C$\subind{4}{2}$}
%			\end{minipage} \\ \vspace{.1cm} 
%			\begin{minipage}{0.5\textwidth}
%				\includegraphics[width=.45\textwidth]{C4_3}
%		%		\subcaption{C$\subind{4}{3}$}
%				\label{fig:C_4^3}
%			\end{minipage}
%			\begin{minipage}{0.4\textwidth}
%				\includegraphics[width=.45\textwidth]{C4_4}
%				\vspace{.8cm}
%		%		\subcaption{C$\subind{4}{4}$}
%			\end{minipage} %\\ \noindent\rule{1.05\textwidth}{0.3pt} \\ \vspace{.1cm} 
%			\begin{minipage}{0.5\textwidth}
%				\includegraphics[width=3.8cm]{C8_1}
%				\vspace{1.15cm}
%		%		\subcaption{C$\subind{8}{1}$}
%			\end{minipage}
%			\begin{minipage}{0.4\textwidth}
%				\includegraphics[width=3.8cm]{C8_2}
%		%		\subcaption{C$\subind{8}{2}$}
%				\label{fig:C_8^2} 
%			\end{minipage} \\ \vspace{-.5cm} \noindent\rule{1.05\textwidth}{0.3pt} \\ \vspace{.1cm} 
%			\begin{minipage}{0.4\textwidth}
%			\centering
%				\includegraphics[width=.6cm]{C16}
%		%		\subcaption{C$\subind{16}{}$}
%			\end{minipage}
%		\end{figure} 
		}
		\end{column}
	\end{columns}
	\only<2>{
	\begin{itemize}
		\item Regla $2^k$
		\item Clases de equivalencia: swaps de partículas y permutaciones de base local
	\end{itemize}
	}
\end{frame}


\begin{frame}[t]{Discusión}
\begin{columns}[t]
	\begin{column}{0.5\textwidth}
	Clasificación:
	\begin{itemize}
		\item Swap de partículas 
		\item Permutaciones de elmentos de base local
	\end{itemize}
	\end{column}
	\begin{column}{0.5\textwidth}
		Reglas:
		\begin{itemize}
			\item Regla $2^k$
			\item Regla espejo
		\end{itemize}			
	\end{column}
\end{columns}
\end{frame}

\section{Canales de Pauli constantes sobre los ejes}
\begin{frame}{Mapeos de Ruskai}
	\begin{itemize}
		\item Qué es un mapa de Ruskai, hablar de MUB en fundamentos teóricos
		o en algún lado?
		\item Los canales PCE no son mapas de Ruskai
	\end{itemize}
\end{frame}

\begin{frame}{Conclusiones}
	\begin{itemize}
		\item Conclusión del trabajo
		\item Trabajo futuro
	\end{itemize}
	
%	\onslide<2>{
	\vspace{1cm}\hfill ¡Gracias!
%	}
\end{frame}

\end{document}
